%% SIGCSE TS 2027 submission --- ACM SIG Conference Proceedings format.
%%
%% Per the SIGCSE TS 2027 CFP: 2-column ACM "sigconf" format, US Letter,
%% double-anonymized. Body limited to 6 pages; references may use a 7th page.
%% Do NOT alter fonts, margins, or the template.
%%
%% NOTE (2027-specific): the CFP says do NOT use the [anonymous] class option.
%% Instead, keep visible author blocks under the title with all identifying
%% information removed, reserving the space (4 lines/author, 3 blocks/row).
%% So we compile with [review] only and anonymize the \author blocks by hand.
%% For camera-ready: remove "review", fill in the real \author blocks, and
%% restore the \acmConference / copyright / DOI block.
%%
%% Build:  make        (pdflatex -> bibtex -> pdflatex x2)
%% acmart.cls and ACM-Reference-Format.bst are vendored in this directory.

\documentclass[sigconf,review]{acmart}

%% --- Anonymous review: suppress the ACM reference + copyright strip ---
\settopmatter{printacmref=false}
\setcopyright{none}
\renewcommand\footnotetextcopyrightpermission[1]{}
\pagestyle{plain}

%% --- Conference / DOI metadata (placeholders until acceptance) ---
\acmConference[SIGCSE TS '27]{Proceedings of the 58th ACM Technical Symposium
  on Computer Science Education}{February 17--20, 2027}{Sacramento, CA, USA}
%% \acmDOI{}        % filled in for camera-ready
%% \acmISBN{}       % filled in for camera-ready

\usepackage{booktabs}
\usepackage{balance}

\begin{document}

%% ===========================================================================
\title{Automated Curriculum Alignment Checking with Coordinated Knowledge Graphs}
%\subtitle{A Tool for AI-Assisted Backwards Design}
%% TODO: title is a working draft --- revisit once the evaluation lands.

%% --- Authors ---
%% Double-anonymized: identifying info removed, but the blocks stay visible to
%% reserve space (CFP: 4 lines/author, up to 3 author blocks per row). Replace
%% with real author/affiliation/email blocks for camera-ready.
\author{Anonymous Author}
\affiliation{%
  \institution{Anonymous Institution}
  \city{Anonymized}
  \country{}
}
\email{anonymous1@example.com}

\author{Anonymous Author}
\affiliation{%
  \institution{Anonymous Institution}
  \city{Anonymized}
  \country{}
}
\email{anonymous2@example.com}

\renewcommand{\shortauthors}{Anonymous et al.}

%% ===========================================================================
\begin{abstract}
Instructors design courses carefully, yet structural misalignments — objectives never assessed, formative practice gaps below a mastery threshold, methods that don't maximize
student learning — are difficult to detect by inspection alone. We present a tool that aligns
three knowledge graphs:  one for the course content, one for the pedagogy and the science of learning, and the third for the course design.  A number of checks are run on the course
design ensuring appropriate scaffolding and structure in topical coverage, student activities and Bloom's taxonomy.  (HELP don't love this)  The tool provides
actionable, specific, research-based recommendations for improving student learning.

We report on a preliminary study for courses throughout the CS curriculum evaluating the recommendations for validity, novelty
and actionability.  In particular, we show that the tool consistently identified at least two actionable findings per course that the instructors indicate were readily adoptable.




Instructor verification is decoupled from representation granularity: instructors confirm fidelity via a rendered weekly outline and a roughly 20-item structural census rather than inspecting individual graph nodes, reducing verification time to approximately 20 minutes.

We report a formative multi-instructor field study (N=4–5) in which each instructor contributed one course they taught, pre-registered their top known issues before seeing tool output, then rated the tool's top findings on validity, novelty, and actionability. Our headline metric is novel-and-actionable findings — issues the tool surfaced that the instructor had not previously identified and would act on. The tool consistently identified at least one novel actionable finding per course. We reflect on encoding fidelity, finding consolidation, failure modes, and practical guidance for instructors considering adoption. All schema definitions, encoding prompts, and check implementations are publicly released.
	
	
%% TODO: tighten once results are in. Per the study design, the abstract must
%% NOT promise a level of "novel" findings the formative data may not deliver.
Backwards design asks instructors to align learning outcomes, activities, and
assessments, but the alignment lives implicitly across a syllabus, slide decks,
and an assessment plan, where gaps are hard to see. We present a tool that makes
this structure explicit and checkable. The tool represents a course as a
\emph{course-design graph} that references two shared upstream knowledge
graphs---one for learning science and pedagogy, one for the course's
content---and runs a suite of automated checks for alignment, coverage,
cognitive-demand progression, scaffolding, and running-example consistency. An
LLM-assisted pipeline encodes a course directly from the materials an instructor
already has (objectives, slide decks, an assessment plan), so adoption does not
require hand-building a graph. Each finding is paired with a concrete, and where
possible closed-loop--validated, remediation. We describe the schema, the
checks, the encoding pipeline, and a worked encoding of a real undergraduate AI
course (a \emph{NNN}-node design graph), and report a formative multi-instructor
study of whether the tool's findings are judged valid, novel, and actionable for
instructors' own courses. \textbf{[Results: TODO.]} The schema, check
definitions, and encoding prompts are released to support adoption.
\end{abstract}

%% --- CCS concepts (pick the final set from the ACM CCS tool) ---
\begin{CCSXML}
<ccs2012>
   <concept>
       <concept_id>10003456.10003457.10003527.10003531</concept_id>
       <concept_desc>Social and professional topics~Computer science education</concept_desc>
       <concept_significance>500</concept_significance>
   </concept>
   <concept>
       <concept_id>10003456.10003457.10003527.10003528</concept_id>
       <concept_desc>Social and professional topics~Curriculum development</concept_desc>
       <concept_significance>300</concept_significance>
   </concept>
   <concept>
       <concept_id>10010147.10010178.10010179</concept_id>
       <concept_desc>Computing methodologies~Knowledge representation and reasoning</concept_desc>
       <concept_significance>100</concept_significance>
   </concept>
</ccs2012>
\end{CCSXML}

\ccsdesc[500]{Social and professional topics~Computer science education}
\ccsdesc[300]{Social and professional topics~Curriculum development}
\ccsdesc[100]{Computing methodologies~Knowledge representation and reasoning}

\keywords{course design, backwards design, constructive alignment, knowledge
graphs, curriculum tools, large language models, instructional design}

\maketitle

%% ===========================================================================
\section{Introduction}
\label{sec:intro}

Backwards design~\cite{wiggins2005understanding,tyler1949basic} and constructive
alignment~\cite{biggs1996constructive} give a clear prescription: start from the
learning outcomes, choose assessments that provide evidence those outcomes were
met, and choose activities that prepare students for those assessments. The
prescription is widely taught and widely endorsed. Yet the artifacts that make up
an actual course---a syllabus, a sequence of slide decks, an assessment plan, a
list of in-class activities---encode this alignment only implicitly and
redundantly. Whether every outcome is actually assessed, whether an assessment
sits at the cognitive level its outcome demands, whether a topic is assessed
before it is taught, or whether a worked example is reused consistently enough to
support transfer, are all questions an instructor must answer by reading the
whole course and holding it in their head at once. They are exactly the questions
that go unanswered when a course is revised one deck at a time.

% --- The gap / why now ---
% TODO: 1 paragraph. Position against (a) curriculum ontologies and KG
% representations of curriculum [christou2025ontology, chen2018knowedu,
% melillan2023software] --- which represent but do not *check* against design
% principles --- and (b) LLM generation of course artifacts
% [hu2024teaching, zheng2025lessonplanlm] --- which generate but are not grounded
% in a checkable, instructor-owned representation. The gap: a tool that makes a
% real course's design *checkable* against backwards-design principles, encoded
% automatically from existing materials.

% --- This paper / contributions ---
We present a tool that makes a course's design explicit and automatically
checkable. The central representation is a \emph{course-design graph} that
references two shared, reusable upstream knowledge graphs---one for learning
science and pedagogy, one for the course's subject-matter content---with all
cross-graph edges pointing outward from the design graph
(\S\ref{sec:overview}). Over this representation the tool runs a suite of
backwards-design checks (\S\ref{sec:checks}), encodes a course from its existing
raw materials via an LLM-assisted pipeline (\S\ref{sec:pipeline}), and turns each
finding into a concrete remediation (\S\ref{sec:recommend}).

This paper contributes:
\begin{itemize}
  \item A \textbf{three-graph representation} for course design that separates
    reusable pedagogy and content knowledge from a course-specific,
    term-agnostic design graph (\S\ref{sec:overview}--\S\ref{sec:schema}).
  \item A suite of \textbf{automated backwards-design checks}---alignment,
    coverage, cognitive-demand progression, scaffolding/prerequisite ordering,
    running-example consistency, and class-day structure---grounded in the
    instructional-design and learning-science literature
    (\S\ref{sec:checks}).
  \item An \textbf{LLM-assisted encoding pipeline} that builds the graph from
    materials instructors already have, lowering the adoption cost
    (\S\ref{sec:pipeline}).
  \item A \textbf{formative, multi-instructor evaluation} of whether the
    findings are valid, novel, and actionable on instructors' own courses, with
    honest reflection on which checks pay off (\S\ref{sec:eval}).
\end{itemize}

To support adoption, we release the schema, the check definitions, and the
encoding prompts.\footnote{Artifact link withheld for anonymous review.}

%% ===========================================================================
\section{Background and Related Work}
\label{sec:related}

\subsection{Backwards design and constructive alignment}
% TODO: Tyler rationale [tyler1949basic] -> Understanding by Design
% [wiggins2005understanding] -> constructive alignment [biggs1996constructive];
% 4C/ID [merrienboer1992training] as the closest full ID framework; Bloom's
% revised taxonomy [anderson2001taxonomy] for the cognitive-demand axis;
% Porter's content x cognitive-demand matrix [porter2002measuring] as the
% established way to quantify coverage. These define the principles the checks
% operationalize.

\subsection{Knowledge graphs and ontologies for curriculum}
% TODO: curriculum/learning-material ontologies [christou2025ontology],
% education KG construction [chen2018knowedu], and the systematic mapping of
% curriculum-design tool support [melillan2023software] (ontologies the most
% common representation). Contrast: these *represent* curriculum; we *check* a
% design against backwards-design principles and emit actionable findings.

\subsection{LLMs for course-artifact generation}
% TODO: GPT-4 teaching-plan generation + automated evaluation
% [hu2024teaching]; KB-grounded lesson-plan generation [zheng2025lessonplanlm].
% Contrast: these generate artifacts; our pipeline encodes an *existing* course
% into an instructor-owned, checkable representation, and our LLM stages are a
% seam inside an otherwise deterministic pipeline.

\subsection{Evidence base for the example checks}
% TODO: worked-examples synthesis [atkinson2000learning]; schema induction and
% analogical transfer [gick1983schema]; analogical encoding via side-by-side
% comparison [gentner2003learning]. These ground checks #13--17 (a running
% example must be present, bridge concepts, and stay structurally consistent).

%% ===========================================================================
\section{System Overview: Three Coordinated Graphs}
\label{sec:overview}

% TODO: figure --- the three graphs and the one-way reference direction.
% \begin{figure}[t] \centering \includegraphics ... \caption{} \end{figure}

% The three graphs and their substrates:
%   - Learning-science graph: how learning works, pedagogy, techniques. Shared.
%   - Content graph: the course's subject material; supports supersession. Per-course.
%   - (AI-education graph: teaching *moves*, grounds the class-day check.) Shared.
%   - Course-design graph: activities, lectures, readings, assessments. Custom,
%     markdown + YAML frontmatter; term-agnostic, versioned via git.
% Reference direction is ONE-WAY: design references the upstream graphs; no
% reverse links. An activity *instantiates* a technique; an activity *covers* a
% content concept. This keeps the shared graphs reusable across courses.

%% ===========================================================================
\section{The Course-Design Schema}
\label{sec:schema}

% TODO: node + edge types; direction families (enablement / reference);
% the "three Blooms" (intended on targets/CLOs, assessed on assessments,
% afforded by activities); grading schemes (weighted/points/mastery/specs);
% thematic (module) vs temporal (session/schedule) axes decoupled.
% Keep this tight --- one labeled example fragment beats a full type catalog.

%% ===========================================================================
\section{Backwards-Design Checks}
\label{sec:checks}

% TODO: organize the ~25 checks into families and report what each protects:
%   - structural alignment (outcome<->assessment<->activity)
%   - coverage depth & sufficiency of evidence; Bloom mismatch
%   - cognitive-demand progression: demand cliffs, premature integration,
%     premature assessment; prerequisite ordering
%   - running examples: presence, concept-bridging, structural consistency (#13-17)
%   - grading-scheme and content-coverage checks
%   - class-day structure / active-learning balance (#25)
% Present as a table (check family | what it catches | grounding citation).
% Note the per-target coverage scorecard.

%% ===========================================================================
\section{Encoding Pipeline}
\label{sec:pipeline}

% TODO: raw materials -> common IR -> course-design graph.
%   - ingest adapters (LaTeX/Beamer, PPTX, Word, PDF; platform adapters)
%   - deterministic stages (schedule order, grading-scheme detection,
%     node/edge extraction) vs LLM stages (deck-outline interpretation,
%     semantic develops/assesses, intended+afforded Bloom) behind a model seam
%   - emphasize: instructor points the tool at existing materials; no hand-built
%     graph. This is part of the contribution, not plumbing.

%% ===========================================================================
\section{From Findings to Fixes}
\label{sec:recommend}

% TODO: consolidation (collapse near-duplicate hits to one root-cause finding),
% impact ranking (= #targets touched x severity weight), and remediation:
% every finding -> a concrete fix, closed-loop validated where auto-applicable
% (apply to scratch copy, re-run checker, keep only if it clears with no new
% fires). Course-specific specifics are agent-authored, not produced by a model
% call inside the scripts --- the tooling stays deterministic.

%% ===========================================================================
\section{Worked Example}
\label{sec:example}

% TODO: the real undergraduate Intro AI course, encoded as a NNN-node design
% graph. Show 1-2 concrete findings end-to-end (finding -> graph location ->
% proposed fix). Use a rendered-graph excerpt as the figure.

%% ===========================================================================
\section{Formative Evaluation}
\label{sec:eval}

% TODO (study in progress --- DO NOT report results until collected):
% Multi-instructor usefulness study (n=4-5), each contributing one course they
% taught. Per course: pipeline encodes it; instructor verifies the encoding
% (layered fidelity check) and pre-registers their top-3 known issues; the tool
% emits findings; instructor rates the top-10 on validity / novelty /
% actionability / localization, then skims the remainder for mis-ranking.
% Headline = novel-AND-actionable findings. Report per-check-category rates
% (which checks pay off). Quote 3-5 verbatim novel-and-actionable vignettes.
% Frame as a field study (ERT register), not a controlled experiment.

\subsection{Method}
% TODO

\subsection{Results}
% TODO --- placeholder until data are collected; no claims yet.

%% ===========================================================================
\section{Limitations and Caveats for Adopters}
\label{sec:limitations}

% TODO (honest reflection, ERT register, not a generalizability defense):
% construct (usefulness vs. learning outcomes); encoding dependence;
% social-desirability / no decoys this round; small non-random n;
% researcher-in-the-loop; top-10 curation; novelty may be rare by construction.

%% ===========================================================================
\section{Conclusion}
\label{sec:conclusion}

% TODO: 1 short paragraph. Restate the contribution and what an adopting
% instructor or tool-builder should take away (which checks are worth
% implementing). Point to the released artifacts.

%% ===========================================================================
\begin{acks}
% Anonymized for review.
\end{acks}

\balance
\bibliographystyle{ACM-Reference-Format}
\bibliography{refs}

\end{document}
